% !TeX program = xelatex
\documentclass[12pt,a4paper]{extarticle}
\input{../common/preamble}
\newcommand{\docname}{Техническое задание}
\newcommand{\doccode}{\hseDocCodeBase\ ТЗ 01-1}
\newcommand{\doccodelu}{\hseDocCodeBase\ ТЗ 01-1-ЛУ}
\input{../common/commands}
\input{../common/styles}
\begin{document}
\sloppy
\input{../common/title_template}


% ============================================
% АННОТАЦИЯ
% ============================================
\section*{АННОТАЦИЯ}
\addcontentsline{toc}{section}{АННОТАЦИЯ}

Техническое задание~--- основной документ, оговаривающий набор требований и~порядок создания программного продукта, в~соответствии с~которым производится разработка программы, ее тестирование и~приемка.

Настоящее Техническое задание на разработку программы «\hseProjectName» содержит следующие разделы: «Введение», «Основания для разработки», «Назначение разработки», «Требования к~программе», «Требования к~программным документам», «Технико-экономические показатели», «Стадии и~этапы разработки», «Порядок контроля и~приемки» и~приложения\cite{gost19101}.

В разделе «Введение» указано наименование и~краткая характеристика области применения программы.

В разделе «Основания для разработки» указан документ, на основании которого ведется разработка и~наименование темы разработки.

В разделе «Назначение разработки» указано функциональное и~эксплуатационное назначение программного продукта.

Раздел «Требования к~программе» содержит основные требования к~функциональным характеристикам, к~интерфейсу, к~надежности, к~условиям эксплуатации, к~составу и~параметрам технических средств, к~информационной и~программной совместимости, к~маркировке и~упаковке, к~транспортировке и~хранению.

Раздел «Требования к~программным документам» содержит предварительный состав программной документации и~специальные требования к~ней.

Раздел «Технико-экономические показатели» содержит ориентировочную экономическую эффективность, предполагаемую годовую потребность, экономические преимущества разработки программы.

Раздел «Стадии и~этапы разработки» содержит стадии разработки, этапы и~содержание работ.

В разделе «Порядок контроля и~приемки» указаны общие требования к~приемке работы.

\clearpage

% ============================================
% СПИСОК ИСПОЛЬЗУЕМЫХ ГОСТ
% ============================================
\noindent
Настоящий документ разработан в~соответствии с~требованиями:

\begin{enumerate}[label=\arabic*), leftmargin=1.25cm]
\item ГОСТ 19.101-77 Виды программ и~программных документов\cite{gost19101}.
\item ГОСТ 19.102-77 Стадии разработки\cite{gost19102}.
\item ГОСТ 19.103-77 Обозначения программ и~программных документов\cite{gost19103}.
\item ГОСТ 19.104-78 Основные надписи\cite{gost19104}.
\item ГОСТ 19.105-78 Общие требования к~программным документам\cite{gost19105}.
\item ГОСТ 19.106-78 Требования к~программным документам, выполненным печатным способом\cite{gost19106}.
\item ГОСТ 19.201-78 Техническое задание. Требования к~содержанию и~оформлению\cite{gost19201}.
\end{enumerate}

Изменения к~данному Техническому заданию оформляются согласно ГОСТ\,19.603-78\cite{gost19603}, ГОСТ\,19.604-78\cite{gost19604}.

\clearpage

% ============================================
% СОДЕРЖАНИЕ
% ============================================
\hypersetup{linkcolor=black}
\tableofcontents
\hypersetup{linkcolor=blue}
\thispagestyle{fancy}
\clearpage

% ============================================
% 1. ВВЕДЕНИЕ
% ============================================
\section{ВВЕДЕНИЕ}

\subsection{Наименование программы}

% Наименования подставляются автоматически: полное — из названия проекта,
% английское и краткое — из настроек проекта (поля «Название работы на
% английском» и «Краткое наименование программы»).
Наименование программы~--- «\projectname».

Наименование программы на английском языке~--- «\hseProjectNameEn».

Краткое наименование программы~--- «\hseProjectNameShort».

\subsection{Краткая характеристика области применения}

% Подсказка: 1–2 абзаца о том, для чего и в какой предметной области
% применяется программа, какие задачи она решает.
Программа «\projectname» является программной системой, предназначенной для \hseFill{основное назначение системы}. Область применения системы охватывает \hseFill{область применения системы}. Система предоставляет \hseFill{интерфейсы взаимодействия с системой} в~пределах предусмотренных функций.

\clearpage

% ============================================
% 2. ОСНОВАНИЯ ДЛЯ РАЗРАБОТКИ
% ============================================
\section{ОСНОВАНИЯ ДЛЯ РАЗРАБОТКИ}

\subsection{Документ(-ы), на основании которого(-ых) ведется разработка}

% Подсказка: сошлитесь на учебный план направления и на утверждённую
% академическим руководителем тему курсового проекта.
Основанием для разработки является учебный план подготовки бакалавров по~направлению \hseFill{код и наименование направления подготовки} и~утверждённая академическим руководителем образовательной программы «\hseSpecialization» тема курсового проекта.

\subsection{Наименование темы разработки}

% Подсказка: приведите наименование темы и её условное обозначение.
Наименование темы разработки~--- «\projectname».

Условное обозначение темы разработки~--- «\hseFill{условное обозначение темы}».

\clearpage

% ============================================
% 3. НАЗНАЧЕНИЕ РАЗРАБОТКИ
% ============================================
\section{НАЗНАЧЕНИЕ РАЗРАБОТКИ}

\subsection{Функциональное назначение}

% Подсказка: перечислите основные функции, которые выполняет программа.
Программа предназначена для выполнения следующих функций:
\begin{itemize}[nosep, leftmargin=1.25cm]
    \item \hseFill{функция 1, например регистрация пользователей};
    \item \hseFill{функция 2, например ведение реестра объектов};
    \item \hseFill{дополните перечень основных функций}.
\end{itemize}

\subsection{Эксплуатационное назначение}

% Подсказка: опишите, кем и в каких условиях программа эксплуатируется.
Эксплуатационное назначение системы состоит в~использовании \hseFill{интерфейсы и сценарии эксплуатации} для \hseFill{целевые задачи пользователей} в~пределах предусмотренных функций.

\clearpage

% ============================================
% 4. ТРЕБОВАНИЯ К ПРОГРАММЕ
% ============================================
\section{ТРЕБОВАНИЯ К ПРОГРАММЕ}

\subsection{Требования к функциональным характеристикам}
\label{sec:functional_reqs}

\subsubsection{Состав выполняемых функций}

% Подсказка: заполните таблицу функциональных требований ниже; сохраняйте
% единый префикс идентификаторов (ТЗ-Ф-01, …) для трассировки в ВКР/ПМИ.
\input{requirements_table}

\clearpage

Функциональные требования к~системе сформированы на~основе анализа предметной области и~целевых сценариев работы пользователей с~системой.

\subsubsection{Организация входных данных}
\label{sec:input_data}

% Подсказка: опишите форматы, протоколы, кодировку, ограничения и правила
% валидации входных данных.
\begin{hseExample}
Внешние запросы клиентов к~публичному API должны выполняться по~протоколу HTTP(S) с~передачей данных в~формате JSON (кодировка UTF-8).

Для всех интерфейсов ввода должны соблюдаться следующие требования:
\begin{itemize}[nosep, leftmargin=1.25cm]
    \item \textbf{Версионирование}: доступ к~API должен осуществляться через префикс версии в~URL (например, \texttt{/v1/...}) для обеспечения обратной совместимости при обновлении программы;
    \item \textbf{Валидация}: все входящие данные должны проходить проверку на~соответствие схемам (\hseFill{укажите средства валидации}) и~типам данных, указанным в~спецификации интерфейса;
    \item \textbf{Ограничения}: максимальный размер тела HTTP-запроса не~должен превышать \hseFill{укажите лимит, например 1 МБ} (если иное не~оговорено для конкретного метода);
    \item \textbf{Пагинация и фильтрация}: запросы на~получение списков объектов должны поддерживать параметры пагинации (\texttt{limit}, \texttt{offset}), сортировки и~фильтрации по~основным атрибутам;
    \item \textbf{Идемпотентность}: операции изменения состояния ресурсов (POST, PUT, DELETE) должны поддерживать механизмы обеспечения идемпотентности при повторных запросах через заголовок \texttt{Idempotency-Key}.
\end{itemize}

% Подсказка: при наличии межсервисного взаимодействия опишите его протоколы
% (например, gRPC, REST, асинхронное событийное взаимодействие).
Межсервисное взаимодействие должно выполняться по~\hseFill{укажите протоколы межсервисного взаимодействия}.

Для трассировки и~корреляции запросов система должна поддерживать сквозной идентификатор запроса (\texttt{X-Request-ID}) в~заголовках HTTP-запросов. Значение идентификатора должно возвращаться в~теле ответов об~ошибках для упрощения диагностики.
\end{hseExample}

\subsubsection{Организация выходных данных}
\label{sec:output_data}

% Подсказка: опишите форматы ответов, коды состояния и структуру сообщений
% об ошибках.
\begin{hseExample}
Система должна обеспечивать формирование ответов на запросы клиентов в~соответствии со спецификациями интерфейсов.

Требования к выходным данным:
\begin{itemize}[nosep, leftmargin=1.25cm]
    \item \textbf{Статус ответа}: для HTTP-запросов должен использоваться стандартный набор кодов состояния (2xx, 4xx, 5xx). При возникновении ошибок сервер должен возвращать HTTP-статус ошибки и~тело ответа в~формате JSON;
    \item \textbf{Гарантия ответа}: сервер должен обеспечивать обработку запроса в~установленные временные лимиты (см. раздел~\ref{sec:temporal_reqs}). При невозможности завершения операции (таймаут, отказ зависимости) система должна возвращать корректный код ошибки (например, 503 Service Unavailable или 504 Gateway Timeout), не~допуская неограниченного «зависания» соединения;
    \item \textbf{Формат данных}: основным форматом выходных данных является JSON (UTF-8).
\end{itemize}

Для HTTP API при возникновении ошибки тело ответа должно передаваться в~виде JSON-объекта следующей структуры:
\begin{itemize}[nosep, leftmargin=1.25cm]
    \item \texttt{error}: объект с~полями \texttt{code} (строковый код ошибки), \texttt{messages} (список сообщений), опционально \texttt{params} (параметры ошибки) и~\texttt{http\_status\_code};
    \item при ошибках валидации (HTTP 422) дополнительно возвращается детальное описание нарушенных ограничений схемы данных.
\end{itemize}
\end{hseExample}

\subsubsection{Требования к временным характеристикам}
\label{sec:temporal_reqs}

% Подсказка: задайте требования к времени отклика, пропускной способности
% и допустимой нагрузке; при необходимости разделите целевой (промышленный)
% и прототипный (стенд разработчика) профили.
\begin{hseExample}
Временные характеристики задаются раздельно для двух режимов эксплуатации: целевого (промышленная конфигурация) и~прототипного (одиночный стенд разработчика). Подтверждение целевых характеристик в~полном объеме требует выделенной инфраструктуры и~выходит за~рамки разработки прототипа; в~рамках настоящих испытаний подтверждается прототипный профиль, заявленный ниже.

\textbf{Целевой профиль} (промышленная конфигурация на~рекомендуемых вычислительных мощностях, см. раздел~\ref{sec:server_hardware}):
\begin{itemize}[leftmargin=1.25cm]
    \item Время обработки синхронных запросов не~должно превышать \hseFill{укажите время, например 500 мс} для 95\% запросов (p95);
    \item При выполнении длительных асинхронных операций время первичного отклика сервера не~должно превышать \hseFill{укажите время отклика};
    \item Указанные показатели должны сохраняться при штатной нагрузке до~\hseFill{число активных сессий} активных пользовательских сессий одновременно;
    \item Целевая интенсивность запросов (RPS) составляет \hseFill{укажите RPS и обоснование расчета}.
\end{itemize}

\textbf{Прототипный профиль} (стенд разработчика):
\begin{itemize}[leftmargin=1.25cm]
    \item Сохраняемая интенсивность запросов составляет не~менее \hseFill{укажите RPS прототипа} при~доле успешных ответов не~менее \hseFill{укажите долю, например 98\%} после периода прогрева;
    \item Время отклика на~прототипном стенде не~нормируется и~определяется ресурсами одной машины; фактические значения p50/p95/p99 фиксируются по~результатам прогона как метрика для сравнения релизов;
    \item Прототипный профиль является условием приемки в~рамках настоящих испытаний и~подтверждается процедурой нагрузочного тестирования (см. соответствующий раздел программы и~методики испытаний).
\end{itemize}
\end{hseExample}

\subsection{Требования к интерфейсу}
\label{sec:interface_reqs}

% Подсказка: общая характеристика интерфейсов взаимодействия с программой.
% Если у программы есть графический интерфейс, добавьте подраздел
% с требованиями к нему.
\begin{hseExample}
Система должна предоставлять программный интерфейс (API) для автоматизации и~интеграций.
\end{hseExample}

\subsubsection{Требования к программному интерфейсу (API)}

% Подсказка: опишите стиль API, способ документирования, авторизацию
% и правила единообразия.
\begin{hseExample}
Программный интерфейс должен обеспечивать полный доступ к~функциям системы для внешних потребителей.
\begin{itemize}[nosep, leftmargin=1.25cm]
    \item \textbf{RESTful API}: внешние методы управления должны быть реализованы в~соответствии с~принципами архитектурного стиля REST (использование HTTP-методов GET, POST, PUT, DELETE);
    \item \textbf{Документирование}: спецификация API должна быть доступна в~автоматизированном режиме в~формате OpenAPI для упрощения интеграции и~тестирования;
    \item \textbf{Авторизация}: доступ к~защищенным методам API должен осуществляться на основе заголовка \texttt{Authorization} с~использованием \hseFill{механизм авторизации, например Bearer-токена JWT};
    \item \textbf{Единообразие}: структура ответов, форматы дат (ISO 8601) и~способы именования полей должны быть унифицированы для всех компонентов программы.
\end{itemize}
\end{hseExample}

\subsection{Требования к надежности}

\subsubsection{Требования к обеспечению надежного (устойчивого) функционирования программы}

% Подсказка: укажите целевые показатели доступности и механизмы устойчивости.
\begin{hseExample}
Система должна сохранять работоспособность и~обеспечивать корректную обработку данных в~условиях штатной эксплуатации. Для обеспечения устойчивого функционирования устанавливаются следующие требования:
\begin{itemize}[nosep, leftmargin=1.25cm]
    \item \textbf{Доступность}: целевой показатель доступности системы (Uptime) должен составлять не~менее \hseFill{укажите показатель, например 99.5\%} (без учета регламентных работ на~инфраструктуре);
    \item \textbf{Самовосстановление}: программные компоненты должны поддерживать механизмы автоматического перезапуска при критических ошибках (\hseFill{укажите механизмы самовосстановления});
    \item \textbf{Изоляция отказов}: отказ одного из~компонентов не~должен приводить к~полной неработоспособности системы (Graceful Degradation); смежные компоненты должны возвращать корректные коды ошибок (например, 503) при недоступности зависимостей;
    \item \textbf{Целостность данных}: при аварийном завершении работы компонентов должна обеспечиваться сохранность данных в~персистентных хранилищах.
\end{itemize}
\end{hseExample}

\subsubsection{Контроль входной информации}

\begin{hseExample}
Входная информация должна соответствовать ограничениям, изложенным в~разделе~\ref{sec:input_data}. В~случае некорректного ввода система должна возвращать соответствующий код ответа (например, 400 Bad Request) и~структурированное описание ошибок валидации.
\end{hseExample}

\subsubsection{Контроль выходной информации}

\begin{hseExample}
Выходная информация должна соответствовать спецификациям, изложенным в~разделе~\ref{sec:output_data}. В~случае невозможности формирования корректного ответа система должна возвращать ошибку с~соответствующим HTTP-статусом (5xx).
\end{hseExample}

\subsubsection{Время восстановления после отказа}

% Подсказка: задайте целевые показатели времени восстановления (RTO/RPO)
% для актуальных в вашем проекте классов инцидентов.
\begin{hseExample}
Для различных классов инцидентов устанавливаются следующие целевые показатели времени восстановления (RTO):
\begin{itemize}[nosep, leftmargin=1.25cm]
    \item \textbf{Отказ программного компонента}: не~более \hseFill{укажите время, например 1 минута} (автоматический перезапуск);
    \item \textbf{Сбой инфраструктурных сервисов} (отказ БД, брокера сообщений): не~более \hseFill{укажите время, например 15 минут};
    \item \textbf{Критический сбой} (полная потеря работоспособности): не~более \hseFill{укажите время, например 1 час} (восстановление из~резервной копии или полная переустановка через автоматизированные сценарии).
\end{itemize}

Целевой показатель точки восстановления (RPO) для персистентных данных должен составлять не~более \hseFill{укажите время, например 15 минут}. Подтверждение соблюдения метрик RTO/RPO должно выполняться путем проведения регулярных испытаний на~стенде, имитирующем профиль нагрузки, описанный в~разделе~\ref{sec:temporal_reqs}.
\end{hseExample}

\subsubsection{Отказы из-за некорректных действий оператора}

\begin{hseExample}
Риск отказа программы вследствие некорректных действий оператора должен быть минимизирован за~счет многоуровневой валидации входных данных. Критические операции изменения конфигурации должны логироваться и~требовать явного подтверждения (например, через механизмы идемпотентности и~предупредительные статусы API).
\end{hseExample}

\clearpage

\subsection{Условия эксплуатации}

\subsubsection{Климатические условия эксплуатации}

Прямые требования к~климатическим условиям эксплуатации программного продукта не~предъявляются. Программа должна функционировать в~пределах климатических норм, предусмотренных для технических средств, на~которых она развернута.

\subsubsection{Требования к численности и квалификации персонала}

% Подсказка: опишите роли, численность и требуемую квалификацию персонала.
Для развертывания и~эксплуатации системы необходим как минимум один квалифицированный специалист, владеющий навыками \hseFill{перечислите ключевые навыки специалиста}.

Взаимодействие с~системой предполагает наличие следующих ролей и~уровней квалификации:
\begin{itemize}[nosep, leftmargin=1.25cm]
    \item \textbf{Конечный пользователь}: \hseFill{опишите требования к пользователю};
    \item \textbf{\hseFill{роль сопровождения системы}}: \hseFill{опишите требования к квалификации}.
\end{itemize}

\subsection{Требования к составу и параметрам технических средств}

\subsubsection{Требования к серверному оборудованию}
\label{sec:server_hardware}

% Подсказка: приведите рекомендуемые и минимальные требования к серверу
% (процессор, ОЗУ, диск, сеть).
Рекомендуемые требования к~вычислительным мощностям серверного оборудования, достаточные для запуска всех компонентов системы:

\begin{itemize}[leftmargin=1.25cm]
\item Процессор: \hseFill{число ядер}
\item Оперативная память: \hseFill{объем ОЗУ}
\item Диск: \hseFill{объем диска}
\item Сеть: \hseFill{пропускная способность сети}
\end{itemize}

Минимальные требования для локальной разработки:

\begin{itemize}[leftmargin=1.25cm]
\item Процессор: \hseFill{число ядер, например 4 ядра}
\item Оперативная память: \hseFill{объем, например 8 ГБ}
\item Диск: \hseFill{объем, например 50 ГБ}
\item Канал: \hseFill{скорость, например 50 Мбит/с}
\end{itemize}

\subsubsection{Требования к клиентскому оборудованию}

Общие требования к~клиентскому оборудованию для работы приложения:

\begin{itemize}[leftmargin=1.25cm]
\item Доступ в~интернет по~протоколу HTTP/HTTPS;
\item Пропускная способность канала не~ниже \hseFill{укажите скорость, например 10 Мбит/с};
\item Задержка (RTT) до~сервера не~более \hseFill{укажите задержку, например 500 мс}.
\end{itemize}

\subsubsection{Требования к инфраструктуре развертывания}
\label{sec:deployment_infra}

% Подсказка: опишите среду развёртывания (например, Docker Compose или
% Kubernetes), функциональные контуры и требования к инфраструктурным
% компонентам: хранилища, брокеры, наблюдаемость, CI/CD.
Система должна быть развернута в~\hseFill{укажите среду развертывания}. Для обеспечения надежности и~изоляции ресурсов эксплуатация должна выполняться в~рамках выделенных функциональных контуров (например, контуры тестирования и~промышленной эксплуатации). Развертывание должно выполняться с~использованием единого стандарта деплоя и~конфигурации.

\textbf{Требования к хранилищам данных и брокерам сообщений:}

\begin{itemize}[leftmargin=1.25cm]
\item \hseFill{укажите СУБД и режим работы};
\item \hseFill{укажите прочие хранилища и брокеры}.
\end{itemize}

\textbf{Требования к наблюдаемости (Observability):}

\begin{itemize}[leftmargin=1.25cm]
\item Сбор и~хранение метрик должны осуществляться в~специализированном хранилище с~визуализацией через графические панели (\hseFill{укажите инструменты мониторинга});
\item Логирование должно быть централизованным и~поддерживать поиск (\hseFill{укажите инструменты логирования}).
\end{itemize}

\textbf{Требования к CI/CD и процессам сборки:}

\begin{itemize}[leftmargin=1.25cm]
\item Сборка и~доставка приложения должны быть автоматизированы в~рамках CI/CD конвейеров (\hseFill{укажите CI/CD-платформу});
\item Должна быть обеспечена автоматическая проверка качества кода (статический анализ).
\end{itemize}

\subsection{Требования к информационной и программной совместимости}

% Подсказка: укажите языки, фреймворки, СУБД, протоколы и среду выполнения.
\begin{itemize}[leftmargin=1.25cm]
\item Архитектура серверной части: \hseFill{укажите архитектурный стиль}.
\item Разрешенные языки программирования: \hseFill{укажите языки и версии}.
\item Предпочитаемые фреймворки и~библиотеки: \hseFill{укажите фреймворки и библиотеки}.
\item Разрешенные хранилища и~брокеры сообщений: \hseFill{укажите СУБД и брокеры}.
\item Среда выполнения и~деплой: \hseFill{укажите среду выполнения и деплоя}; конфигурация через переменные окружения; секреты не~должны храниться в~исходном коде.
\end{itemize}

\subsection{Требования к маркировке и упаковке}

\subsubsection{Требования к упаковке и составу поставки}

% Подсказка: опишите состав дистрибутива и формат поставки.
\begin{hseExample}
Программный продукт должен поставляться в~виде цифрового дистрибутива, включающего следующие компоненты:
\begin{itemize}[nosep, leftmargin=1.25cm]
    \item \textbf{Дистрибутив программы}: \hseFill{укажите форму поставки, например Docker-образы};
    \item \textbf{Конфигурации развертывания}: \hseFill{укажите артефакты развертывания};
    \item \textbf{Манифест релиза}: файл (например, \texttt{release.json} или \texttt{RELEASENOTES.md}), содержащий перечень версий всех компонентов, входящих в~поставку;
    \item \textbf{Программная документация}: комплект документов в~электронном виде согласно разделу~\ref{sec:doc_composition}.
\end{itemize}
\end{hseExample}

\subsubsection{Требования к маркировке и версионированию}

\begin{hseExample}
Для всех компонентов системы устанавливаются следующие правила маркировки:
\begin{itemize}[nosep, leftmargin=1.25cm]
    \item \textbf{Версионирование программного кода}: должно соответствовать спецификации Semantic Versioning 2.0.0 (формат \texttt{MAJOR.MINOR.PATCH});
    \item \textbf{Маркировка артефактов поставки}: каждый артефакт должен быть помечен тегом, соответствующим версии релиза;
    \item \textbf{Метаданные сборки}: артефакты должны содержать сведения о~дате сборки, идентификаторе коммита (Git SHA) и~ответственном за~сборку;
    \item \textbf{Контрольные суммы}: для каждого артефакта поставки должна быть сформирована контрольная сумма (SHA-256) для проверки целостности при получении дистрибутива.
\end{itemize}
\end{hseExample}

\subsection{Требования к транспортированию и хранению}

\subsubsection{Требования к хранению}

\begin{hseExample}
Хранение всех составляющих программного продукта должно осуществляться в~цифровом виде в~защищенных хранилищах:
\begin{itemize}[nosep, leftmargin=1.25cm]
    \item \textbf{Исходный код и конфигурации}: должны храниться в~системе контроля версий (Git) с~настроенным разграничением прав доступа;
    \item \textbf{Артефакты сборки}: должны храниться в~реестре артефактов (\hseFill{укажите реестр артефактов}), поддерживающем версионирование;
    \item \textbf{Секреты}: пароли, ключи и~сертификаты должны храниться отдельно от~программного кода в~специализированном хранилище.
\end{itemize}
\end{hseExample}

\subsubsection{Требования к транспортированию}

\begin{hseExample}
Транспортирование (доставка) программного продукта должно выполняться через защищенные каналы связи:
\begin{itemize}[nosep, leftmargin=1.25cm]
    \item \textbf{Механизм доставки}: передача артефактов между средами должна осуществляться автоматизированными CI/CD конвейерами;
    \item \textbf{Защита каналов}: все операции по~транспортированию артефактов и~конфигураций должны выполняться по~протоколам HTTPS или SSH;
    \item \textbf{Проверка целостности}: при автоматизированном развертывании система должна выполнять проверку контрольных сумм и~подписей артефактов (если применимо) перед запуском в~целевой среде;
    \item \textbf{Изоляция сред}: прямая передача данных между средами разработки и~промышленной эксплуатации должна быть исключена; доставка осуществляется через промежуточный реестр артефактов.
\end{itemize}
\end{hseExample}

\clearpage

% ============================================
% 5. ТРЕБОВАНИЯ К ПРОГРАММНОЙ ДОКУМЕНТАЦИИ
% ============================================
\section{ТРЕБОВАНИЯ К ПРОГРАММНОЙ ДОКУМЕНТАЦИИ}

\subsection{Состав программной документации}
\label{sec:doc_composition}

% Подсказка: перечислите состав программной документации с указанием ГОСТов —
% оставьте только те руководства, которые входят в ваш комплект
% (программиста и/или оператора).
\begin{enumerate}[label=\arabic*), leftmargin=1.25cm]
\item «\projectname». Техническое задание (ГОСТ~19.201-78).
\item «\projectname». Пояснительная записка (ГОСТ~19.404-79).
\item «\projectname». Программа и~методика испытаний (ГОСТ~19.301-79).
\item «\projectname». Текст программы (ГОСТ~19.401-78).
\hseManualDocItems
\end{enumerate}

\subsection{Специальные требования к программной документации}
\label{sec:doc_reqs}

Документы к~программе должны быть выполнены в~соответствии с~ГОСТ~19.106-78 и~ГОСТами к~каждому виду документа (см.~раздел~\ref{sec:doc_composition}).

\clearpage

% ============================================
% 6. ТЕХНИКО-ЭКОНОМИЧЕСКИЕ ПОКАЗАТЕЛИ
% ============================================
\section{ТЕХНИКО-ЭКОНОМИЧЕСКИЕ ПОКАЗАТЕЛИ}

\subsection{Ориентировочная экономическая эффективность}

% Подсказка: приведите ориентировочную экономическую эффективность
% или укажите, что она не предусмотрена.
В~рамках данной работы экономическая эффективность не~предусмотрена.

\subsection{Предполагаемая потребность}

% Подсказка: опишите потенциальных пользователей и предполагаемую потребность.
Разрабатываемая система может представлять интерес для \hseFill{целевая аудитория разработки}. Потенциальными пользователями являются \hseFill{перечислите потенциальных пользователей}.

\subsection{Экономические преимущества разработки по~сравнению с~отечественными и~зарубежными аналогами}

% Подсказка: кратко опишите преимущества разработки или сошлитесь на анализ
% аналогов в пояснительной записке.
Сравнительная оценка отечественных и~зарубежных аналогов не~включается в~техническое задание, поскольку данный документ фиксирует назначение, требования, состав поставки и~порядок приемки разрабатываемой программы. Анализ аналогов, критерии сравнения и~вывод об~актуальности разработки приведены в~тексте пояснительной записки.

\clearpage

% ============================================
% 7. СТАДИИ И ЭТАПЫ РАЗРАБОТКИ
% ============================================
\section{СТАДИИ И ЭТАПЫ РАЗРАБОТКИ}

\subsection{Стадии разработки, этапы и содержание работ}

% Подсказка: заполните таблицу стадий, этапов и содержания работ
% (с учётом ГОСТ 19.102-77).
Стадии и~этапы разработки были выявлены с~учетом ГОСТ~19.102-77\cite{gost19102}.

\begin{table}[H]
\caption{Стадии и этапы разработки}
\label{tab:stages}
\small
\begin{tabularx}{\textwidth}{|>{\raggedright\arraybackslash}p{3.5cm}|>{\raggedright\arraybackslash}p{3.5cm}|>{\raggedright\arraybackslash}X|}
\hline
\textbf{Стадии разработки} & \textbf{Этапы работ} & \textbf{Содержание работ} \\
\hline
1. Техническое задание & Подготовительные работы & --- Постановка задачи.\newline
--- Сбор исходных материалов.\newline
--- Предварительный выбор методов решения задач. \\
\hline
& Разработка и утверждение ТЗ & --- Определение требований к программе.\newline
--- Определение требований к техническим средствам.\newline
--- Определение стадий, этапов и сроков разработки программы и документации на нее.\newline
--- Согласование и утверждение ТЗ. \\
\hline
2. Рабочий проект & Разработка программы & --- Разработка и отладка программы. \\
\hline
& Разработка программной документации & --- Разработка программных документов в соответствии с требованиями ГОСТ 19.101-77\cite{gost19101}. \\
\hline
& Испытания программы & --- Разработка, согласование и утверждение порядка и методики испытаний.\newline
--- Корректировка программы и программной документации по результатам испытаний. \\
\hline
3. Внедрение & Подготовка и передача программы & --- Подготовка и передача программы и программной документации для сопровождения. \\
\hline
\end{tabularx}
\end{table}

\subsection{Сроки разработки и исполнители}

% Подсказка: укажите сроки завершения работы и исполнителя(-ей).
Программный продукт (программа и~документация) должен быть завершен не~позднее утвержденного приказом декана ФКН НИУ ВШЭ срока защиты курсового проекта.

Исполнитель~--- \hseAuthorName.

\clearpage

% ============================================
% 8. ПОРЯДОК КОНТРОЛЯ И ПРИЕМКИ
% ============================================
\section{ПОРЯДОК КОНТРОЛЯ И ПРИЕМКИ}

\subsection{Виды испытаний}

Производится проверка корректного выполнения программой заложенных в~нее функций, то~есть осуществляется функциональное тестирование программы. Функциональное тестирование осуществляется в~соответствии с~документом «\projectname». Программа и~методика испытаний (ГОСТ~19.301-79).

\subsection{Общие требования к приемке работы}

Прием программы будет утвержден при корректной работе программы в~соответствии с~разделом~\ref{sec:functional_reqs} данного документа и~при предоставлении полной документации к~продукту, указанной в~разделе~\ref{sec:doc_composition}, выполненной в~соответствии с~требованиями в~разделе~\ref{sec:doc_reqs} данного технического задания.

\clearpage
% Страницы источников печатаются без нижнего штампа (как в реальных
% комплектах) — только номер листа и код документа сверху.
\pagestyle{appendix}

% ============================================
% СПИСОК ИСПОЛЬЗОВАННЫХ ИСТОЧНИКОВ
% ============================================
\section*{СПИСОК ИСПОЛЬЗОВАННЫХ ИСТОЧНИКОВ}
\addcontentsline{toc}{section}{СПИСОК ИСПОЛЬЗОВАННЫХ ИСТОЧНИКОВ}

% Подсказка: добавьте ПЕРЕД ГОСТами источники по технологиям вашего
% проекта в алфавитном порядке, по образцу:
% \bibitem{fastapi}
% FastAPI [Электронный ресурс]. Режим доступа: \url{https://fastapi.tiangolo.com}, свободный. (дата обращения: 14.03.2026).
\renewcommand{\refname}{}
\begin{thebibliography}{99}
\bibitem{gost19101}
ГОСТ 19.101-77: Виды программ и~программных документов. // Единая система программной документации.~--- М.: ИПК Издательство стандартов, 2001.

\bibitem{gost19102}
ГОСТ 19.102-77: Стадии разработки. // Единая система программной документации.~--- М.: ИПК Издательство стандартов, 2001.

\bibitem{gost19103}
ГОСТ 19.103-77: Обозначения программ и~программных документов. // Единая система программной документации.~--- М.: ИПК Издательство стандартов, 2001.

\bibitem{gost19104}
ГОСТ 19.104-78: Основные надписи. // Единая система программной документации.~--- М.: ИПК Издательство стандартов, 2001.

\bibitem{gost19105}
ГОСТ 19.105-78: Общие требования к~программным документам. // Единая система программной документации.~--- М.: ИПК Издательство стандартов, 2001.

\bibitem{gost19106}
ГОСТ 19.106-78: Требования к~программным документам, выполненным печатным способом. // Единая система программной документации.~--- М.: ИПК Издательство стандартов, 2001.

\bibitem{gost19201}
ГОСТ 19.201-78: Техническое задание. Требования к~содержанию и~оформлению. // Единая система программной документации.~--- М.: ИПК Издательство стандартов, 2001.

\bibitem{gost19603}
ГОСТ 19.603-78: Общие правила внесения изменений. // Единая система программной документации.~--- М.: ИПК Издательство стандартов, 2001.

\bibitem{gost19604}
ГОСТ 19.604-78: Правила внесения изменений в~программные документы, выполненные печатным способом. // Единая система программной документации.~--- М.: ИПК Издательство стандартов, 2001.
\end{thebibliography}

\clearpage

% ============================================
% ПРИЛОЖЕНИЕ 1 - ГЛОССАРИЙ
% ============================================
\section*{ПРИЛОЖЕНИЕ 1}
\addcontentsline{toc}{section}{ПРИЛОЖЕНИЕ 1}

\begin{center}
\textbf{ГЛОССАРИЙ}
\end{center}

% Подсказка: добавьте термины предметной области и технологий вашего
% проекта. Порядок — алфавитный: сначала латиница, затем кириллица.
\begin{enumerate}[label=\arabic*., leftmargin=1.25cm]
\item \textbf{\hseFill{термин на латинице}}~--- \hseFill{определение термина}.
\item \textbf{\hseFill{ещё один термин}}~--- \hseFill{определение термина}.
\end{enumerate}

\clearpage


% ============================================
% ЛИСТ РЕГИСТРАЦИИ ИЗМЕНЕНИЙ
% ============================================
\input{../common/change_log}

\end{document}
